\chapter{Introdução}
\label{chapter:introduction}

Teste de software é uma das atividades presentes no processo de desenvolvimento e manutenção de software. De acordo com o \textit{Guide to the Software Engineering Body of Knowledge} (SWEBOK) \cite{bourque:etal:2014}, teste de software consiste na verificação dinâmica de que um programa provê um comportamento esperado a partir de um conjunto finito de casos de teste, adequadamente selecionados a partir de um domínio de execução potencialmente infinito. Por muito tempo, teste de software foi visto sobretudo como uma etapa de verificação da qualidade do software produzido, e não como meio de construir um produto com qualidade \cite{murugesan:1994}. Nesse sentido, testes eram frequentemente realizados apenas ao final da etapa de codificação \cite{yang:etal:2011}, com o objetivo de detectar falhas e fornecer um indicativo de que o software funciona conforme o esperado \cite{bertolino:2007}. Com o passar do tempo, a percepção sobre a utilidade do teste evoluiu. A atividade deixou de ser vista como etapa posterior ao desenvolvimento e passou a estar presente durante todo o processo de desenvolvimento e manutenção de software \cite{bourque:etal:2014}. \citeonline{meszaros:2007} considera a atividade de teste como parte intrínseca do processo de desenvolvimento de software. Assim, a atividade de teste de software passou a incorporar propósitos variados, como detectar falhas \cite{tiwari:etal:2013}, prevenir que erros sejam introduzidos \cite{beizer:1990}, prover um indicativo de que o software funciona \cite{bertolino:2007}, auxiliar na compreensão do projeto e do código do sistema \cite{freeman:etal:2009} e servir como meio de especificação \cite{alvestad:2007} e de documentação \cite{haugset:etal:2008} dos requisitos.

Nesse contexto, o código de teste deixou de ser um artefato auxiliar e passou a exigir os mesmos cuidados aplicados ao código de produção. Assim como o código da aplicação, o código de teste precisa ser mantido, compreendido e ajustado durante todo o projeto \cite{greiler:etal:2013:b}. Para \citeonline{meszaros:2007}, a atividade de teste não deveria aumentar o esforço total de desenvolvimento e manutenção do software, uma vez que o esforço gasto com essa atividade deve ser compensado pela melhoria da qualidade do software. Entretanto, manter o código de teste apresenta desafios, principalmente quando a codificação e a evolução dos testes são realizadas diretamente pelo desenvolvedor. Conforme o software evolui e a quantidade de código aumenta, torna-se mais difícil manter os testes com consistência e sem perder o controle sobre eles. Um fenômeno semelhante ocorre no código de produção: sem refatoração contínua, cada nova funcionalidade tende a exigir mais esforço do que a anterior \cite{fowler:1999}. Com os testes, observa-se comportamento análogo. Novos testes passam a ser mais difíceis de adicionar e a duplicação tende a aumentar \cite{berner:etal:2005}. Isso pode levar à degradação da manutenibilidade do código de teste \cite{greiler:etal:2013:a}. Segundo \citeonline{emery:2009}, um motivo comum para o abandono de testes automatizados é o fato de os testes se tornarem frágeis e de alto custo de manutenção, o que pode tornar os desenvolvedores reticentes tanto em adicionar novos testes quanto em executar os existentes. Mesmo em um cenário em que tanto o código de produção quanto o código de teste são gerados por modelos de Inteligência Artificial (IA), a manutenibilidade do código de teste é um problema que precisa ser tratado. \citeonline{shamshiri:etal:2018} mostra que testes gerados automaticamente afetam o tempo e a confiança do desenvolvedor, e que as tarefas de manutenção levaram 29\% mais tempo quando realizadas com testes gerados automaticamente. Nesse sentido, independentemente do contexto, a manutenibilidade do código de teste ainda é um aspecto importante no processo de desenvolvimento de software.

Evitar a duplicação no código de teste é uma das formas de melhorar a manutenibilidade. Isso ocorre porque mudanças na aplicação passam a ter impacto menor nos testes, e o desenvolvedor dispõe de menos código para manter e gerenciar \cite{berner:etal:2005,greiler:etal:2013:a}. De acordo com \citeonline{berner:etal:2005}, testes tendem a ser repetitivos e apresentam alto potencial de reuso; conforme o software evolui e o número de testes aumenta, é comum que a quantidade de código duplicado entre os testes também aumente. As Figuras~\ref{code:1-cancel-student-enrollment-test} e~\ref{code:2-enroll-student-twice-test} ilustram duplicação no código de teste. Cada figura apresenta um método de teste em uma classe distinta. Os exemplos utilizam a linguagem Kotlin e o framework JUnit 5\footnote{\url{https://junit.org/}}. O código das linhas 3--8 em ambos os testes configura o sistema em teste (\textit{System Under Test} -- SUT) da mesma forma, o que produz código de configuração duplicado.

\begin{figure}[htb]
	\centering
	\lstinputlisting{code/1-cancel-student-enrollment-test.txt}
	\caption{Método de teste com configuração do SUT no próprio método, suscetível à duplicação de código}
	\label{code:1-cancel-student-enrollment-test}
	\fonte{Elaborado pelo autor.}
\end{figure}

\begin{figure}[htb]
	\centering
	\lstinputlisting{code/2-enroll-student-twice-test.txt}
	\caption{Método de teste com lógica de configuração do SUT idêntica à da Figura~\ref{code:1-cancel-student-enrollment-test}, ilustrando a duplicação}
	\label{code:2-enroll-student-twice-test}
	\fonte{Elaborado pelo autor.}
\end{figure}

Técnicas que reduzam essa duplicação podem contribuir para preservar a manutenibilidade do código de teste ao longo do projeto, bem como para diminuir o esforço total de desenvolvimento dos testes. Uma dessas técnicas é a configuração implícita (ou \textit{implicit setup}), em que os testes são organizados em classes de teste, de forma que o código de configuração comum possa ser centralizado em um método de configuração \cite{meszaros:2007}. A Figura~\ref{code:3-university-tests} ilustra a aplicação da configuração implícita. Além da remoção da duplicação de código, esse exemplo mostra como a utilização da configuração implícita traz mais clareza à compreensão do teste. O contexto comum (ter um aluno matriculado em uma disciplina) aparece apenas uma vez e cada método de teste foca no aspecto mais importante para aquele teste.

\begin{figure}[htb]
	\centering
	\lstinputlisting{code/3-university-tests.txt}
	\caption{Suíte de testes refatorada com configuração implícita, centralizando a lógica comum de inicialização e eliminando a duplicação}
	\label{code:3-university-tests}
	\fonte{Elaborado pelo autor.}
\end{figure}

Entretanto, nem sempre é simples decidir como os testes devem ser organizados. Por exemplo, considere o teste da Figura~\ref{code:4-enroll-student-in-a-closed-course-test}. Esse teste requer um estado do SUT diferente do estado comum dos testes da Figura~\ref{code:3-university-tests}. Em vez de uma disciplina com um aluno matriculado, esse teste requer uma disciplina fechada e sem alunos matriculados. Como nem toda configuração pode ser compartilhada, o novo teste não poderia ser colocado na mesma classe de teste dos demais sem que o código da classe existente fosse adaptado. A pergunta que fica é: compensaria fazer essa refatoração? Por um lado, o código que é comum ao novo teste e aos demais (linhas 3--7 da Figura~\ref{code:4-enroll-student-in-a-closed-course-test}) deixaria de ser duplicado. Por outro lado, o código que realiza a matrícula do aluno na disciplina (linha 12 da Figura~\ref{code:3-university-tests}) precisaria ser duplicado nos dois testes que a exigem. Considerando apenas o reuso de código, é tácito concluir que compensaria colocar o novo teste na mesma classe de teste dos demais, uma vez que o reuso obtido seria maior que a duplicação introduzida; para reusar cinco linhas de código, uma linha de código precisaria ser duplicada duas vezes. Mas, e se ao invés de apenas dois, a classe existente já tivesse seis testes? Nesse caso, mover o novo teste para a mesma classe dos demais testes introduziria mais linhas duplicadas do que aquelas que seriam reusadas.

\begin{figure}[htb]
	\centering
	\lstinputlisting{code/4-enroll-student-in-a-closed-course-test.txt}
	\caption{Método de teste que requer um estado distinto do SUT, impedindo o reuso da configuração comum da Figura~\ref{code:3-university-tests}}
	\label{code:4-enroll-student-in-a-closed-course-test}
	\fonte{Elaborado pelo autor.}
\end{figure}

O problema fica evidente mesmo com um exemplo simples: organizar os métodos de teste levando em conta a redução da duplicação de código não é uma tarefa trivial. Em especial, em projetos com centenas ou milhares de testes, a distribuição de um novo teste leva a uma explosão combinatória de possibilidades. Convém destacar que não é apenas uma decisão sobre qual é a melhor classe de teste para inserir o novo teste, mas também sobre como re-distribuir os testes já existentes entre as classes de teste da maneira mais eficiente. Quando um novo teste é introduzido, por exemplo, é possível que a nova distribuição ótima acabe separando testes que antes pertenciam à mesma classe de teste. Percebe-se, então, que com o objetivo de atingir uma distribuição ótima dos testes considerando a redução da duplicação de código, é necessário que o código de teste seja continuamente refatorado e reorganizado.

\section{Problema de pesquisa}

Diante desse contexto, o problema de pesquisa é enunciado da seguinte maneira: como reduzir a duplicação de código de teste sem aumentar o esforço total de desenvolvimento e manutenção e sem alterar o resultado dos testes?

\section{Hipótese de pesquisa}

Considerando o problema apresentado, a hipótese de pesquisa deste trabalho é enunciada da seguinte forma: é possível reduzir a duplicação de código de teste e o esforço de desenvolvimento por meio de um método de refatoração automática que trate o reuso de código de teste como um problema de otimização global da distribuição de métodos entre classes de teste.

\section{Proposta}

Em linha com a hipótese apresentada, este trabalho propõe a utilização de métricas de similaridade e algoritmos de agrupamento para resolver o problema de otimização global da distribuição de métodos entre classes de teste. As métricas de similaridade permitem identificar testes com código de configuração em comum; os algoritmos de agrupamento, por sua vez, propõem grupos de testes cuja redistribuição maximiza o reuso de código pela estratégia de configuração implícita.

\section{Objetivos}

O objetivo geral deste trabalho consiste em reduzir a duplicação de código de teste e o esforço total de desenvolvimento e manutenção dos testes, sem alterar o resultado dos testes. Para alcançar o objetivo geral, os seguintes objetivos específicos deverão ser alcançados: (1) investigar métricas de similaridade para identificar código de configuração em comum entre testes; (2) investigar algoritmos de agrupamento para propor distribuições de métodos entre classes de teste; (3) propor uma arquitetura em pipeline para refatoração de código de teste; (4) especificar e implementar um framework que instancia essa arquitetura sem alterar o comportamento observado dos testes; e (5) avaliar a abordagem quanto à redução de duplicação e ao esforço de refatoração.

\section{Organização do trabalho}

Os demais capítulos deste trabalho são organizados da seguinte forma: no \autoref{chapter:background} é apresentada a fundamentação teórica; no \autoref{chapter:state-of-the-art} é apresentado o estado da arte e os trabalhos relacionados; no \autoref{chapter:proposal} é apresentada a proposta deste trabalho; e no \autoref{chapter:conclusion} é apresentada a conclusão e os próximos passos para elaboração deste trabalho.
